\Askhsh{1777}{4}
\begin{ylh}{1}
    \item 5.2. Παραλληλόγραμμα
    \item 5.3. Ορθογώνιο
    \item 5.6. Εφαρμογές στα τρίγωνα
    \item 5.8. Το ορθόκεντρο τριγώνου.
\end{ylh}
Δίνονται οξυγώνιο τρίγωνο $\mathrm{A}\mathrm{B}\Gamma$, $\mathrm{B}\mathrm{E}, \Gamma\mathrm{Z}$, τα ύψη από τις κορυφές $\mathrm{B}, \Gamma$ αντίστοιχα και $\mathrm{H}$ το ορθόκεντρο του τριγώνου. Επίσης δίνονται τα $\mathrm{M}, \mathrm{N}, \mathrm{K}, \Lambda$ μέσα των ευθυγράμμων τμημάτων $\mathrm{A}\mathrm{B}, \mathrm{A}\Gamma, \Gamma\mathrm{H}, \mathrm{B}\mathrm{H}$ αντίστοιχα.
\begin{alist}
    \item Να αποδείξετε ότι:
    \begin{rlist}
        \item $\mathrm{M}\mathrm{N} = \Lambda\mathrm{K}$ \monades{6}
        \item $\mathrm{N}\mathrm{K} = \mathrm{M}\Lambda = \frac{\mathrm{A}\mathrm{H}}{2}$ \monades{6}
        \item Το τετράπλευρο $\mathrm{M}\mathrm{N}\mathrm{K}\Lambda$ είναι ορθογώνιο. \monades{6}
    \end{rlist}
    \item Αν το $\mathrm{O}$ είναι το μέσο της $\mathrm{B}\Gamma$, να αποδείξετε ότι $\mathrm{M}\hat{\mathrm{O}}\mathrm{K} = 90^\circ$. \monades{7}
\end{alist}
